\section{Análisis, Trazabilidad y Lecciones}
\label{sec:analisis_trazabilidad}

\subsection{Matriz de Trazabilidad Actualizada}

Reflejo de la cobertura basada en los estados reales de aprobación tras ejecutar las pruebas.

\begin{xltabular}{\textwidth}{l X l X}
    \caption{Matriz de Trazabilidad}
    \label{tab:matriz_trazabilidad} \\
    \toprule
    \textbf{Requisito} & \textbf{Descripción} & \textbf{Casos asociados} & \textbf{Estado de los Casos} \\
    \midrule
    RF-01 & Registro de incidencias & CP-F-01, CP-F-02, CP-V-03, CP-F-10, CP-F-13 & 5/5 Aprobados \\
    RF-02 & Auditoría automática & CP-F-15 & Fallido (No implementado) \\
    RF-03 & Gestión de estados & CP-F-04, CP-F-05, CP-F-07, CP-F-12, CP-V-04 & 2 Aprobados, 3 Fallidos \\
    RF-04 & Ubicación geográfica & CP-F-11, CP-V-01, CP-V-02 & 3/3 Aprobados \\
    RF-05 & Roles y permisos & CP-F-06, CP-S-02, CP-S-05 & 3/3 Aprobados \\
    RF-06 & Clasificación por prioridad & CP-F-09 & 1/1 Aprobado \\
    RF-07 & Instituciones responsables & CP-F-08 & Fallido (No implementado) \\
    RNF-01 & API REST & CP-F-14, CP-S-01, CP-S-03, CP-S-04 & 3 Aprobados, 1 Fallido \\
    RNF-02 & Eliminación lógica & CP-F-03 & 1/1 Aprobado \\
    RNF-03 & Docker & CP-F-14 & Fallido \\
    RNF-04 & Precisión GPS & CP-V-01, CP-V-02 & 2/2 Aprobados \\
    \bottomrule
\end{xltabular}

\subsection{Análisis de Zonas Frágiles y Conclusiones}

\begin{itemize}
    \item \textbf{Zonas Frágiles Detectadas:}
    \begin{itemize}
        \item La integración con la API de mapas (Leaflet) para geolocalización presentó inconsistencias al no encontrar parroquias específicas según la latitud y longitud, lo que derivó en la apertura del BUG-01.
        \item Existen varios módulos periféricos que aún no han sido desarrollados en el código fuente, particularmente la gestión de comentarios, el historial de auditoría de cambios y las notificaciones automáticas, ocasionando fallos en bloque sobre los requisitos asociados.
    \end{itemize}
    \item \textbf{Conclusiones Predictivas sobre la Calidad:}
    \begin{itemize}
        \item El sistema presenta una robustez destacable en flujos primarios (CRUD de incidencias, validación en base de datos, seguridad JWT y roles). No obstante, el producto todavía no es apto para un pase a producción real dado que los complementos de interacción al usuario (comentarios y notificaciones) aún se encuentran pendientes.
    \end{itemize}
    \item \textbf{Lecciones Aprendidas:}
    \begin{itemize}
        \item Se deben implementar mecanismos de contingencia (fallbacks manuales) al depender de librerías de terceros; por ejemplo, si el mapa falla al buscar ubicaciones, el usuario debe tener control total para sobreescribir la información.
        \item Es recomendable definir con mayor claridad el alcance de cada iteración, para no ejecutar casos de prueba que inevitablemente fallarán por corresponder a funcionalidades programadas para Sprints posteriores.
    \end{itemize}
\end{itemize}
